\section{第二讲：仿射 Springer 纤维}\label{s:ASF}

仿射 Springer 纤维是环路群的 Springer 纤维的类比物，由 Kazhdan 与 Lusztig \cite{KL} 引入。粗略地说，对经典群而言，Springer 纤维分类的是 $k$ 向量空间中被 $k$ 线性变换固定的旗；仿射 Springer 纤维分类的则是 $F$ 向量空间中被 $F$ 线性变换固定的格（及其链），其中 $F=k\lr{t}$。仿射 Springer 纤维的上同调群带有仿射 Weyl 群的作用。

本节的设定与 \S\ref{ss:notation} 相同。

\subsection{环路群、parahoric 子群与仿射旗簇}
设 $F=k\lr{t}$ 为单变量 $t$ 的形式 Laurent 级数域。$F$ 有离散赋值 $\val_{F}:F^{\times}\to\ZZ$，满足 $\val_{F}(t)=1$，其赋值环为 $\calO_{F}=k\tl{t}$。

\subsubsection{环路群} {\em 正环路群} $L^{+}G$ 是 $k$ 代数上的群值函子。对任意 $k$ 代数 $R$，定义 $LG(R):=G(R\tl{t})$。可以证明 $LG$ 可由 $k$ 上的一个概形表示，但该概形不是有限型的。\zy{原文此处称 $L^{+}G$ 为「正环路群」，紧接的定义却是 $LG(R):=G(R\tl{t})$，即环路群。下一段才定义正环路群 $L^{+}G$。原文的这句应属笔误，译文照录，并登记为勘误 E1。［来源：原文 \texttt{ZhiweiYunPCMI.tex} 第 1251--1253 行文脉。］}

例如当 $G=\GL_{n}$ 时，$LG(R)=\GL_{n}(R\tl{t})$ 中的元素由 $n^{2}$ 个形式 Laurent 级数 $a_{ij}=\sum_{s\geq 0}a^{(s)}_{ij}t^{s}$（$1\leq i,j\leq n$）给出，其中 $a^{(s)}_{ij}\in R$，并满足一个条件：$\det((a_{ij}))$（它是 $a^{(s)}_{ij}$ 的 $n$ 次多项式）在 $R$ 中可逆。因此此时 $LG$ 是以 $a^{(s)}_{ij}$（$1\leq i,j\leq n$，$s\geq 0$）为坐标的无穷维仿射空间中的开子概形。

类似地可以定义{\em 环路群} $LG$ 为 $k$ 代数 $R$ 上的函子 $LG(R)=G(R\lr{t})$。函子 $LG$ 不再可由概形表示，而是由{\em 归纳概形}表示。归纳概形是范畴意义下的归纳极限 $\varinjlim_{m} X_{m}$：即 $\{X_{m}\}$ 构成 $k$ 上概形的归纳系，而 $\varinjlim_{m} X_{m}$ 是函子 $R\mapsto \varinjlim_{m} X_{m}(R)$。当 $G=\GL_{n}$ 时，可令 $X_{m}$ 为 $LG$ 的子函子，使 $X_{m}(R)$ 由元素落在 $t^{-m}R\tl{t}\subset R\lr{t}$ 中的 $n$ 阶可逆矩阵组成。此时与 $L^{+}G$ 的情形同样的论证表明 $X_{m}$ 可由 $k$ 上概形表示。当 $m<m'$ 时，有自然的闭嵌入 $X_{m}\incl X_{m'}$，而此时的 $LG$ 就是归纳极限 $\varinjlim_{m}X_{m}$。对一般的 $G$，参见 \cite[\S1]{BL} 与 \cite[\S2]{Faltings}。

\subsubsection{仿射 Grassmann 簇}\label{sss:Gr} $G$ 的仿射 Grassmann 簇 $\Gr_{G}$ 定义为预层 $R\mapsto LG(R)/L^{+}G(R)$ 在 $k$ 代数范畴中关于 fpqc 拓扑的层化。特别地，有 $\Gr_{G}(k)=G(F)/G(\calO_{F})$。

当 $G=\GL_{n}$ 时，仿射 Grassmann 簇 $\Gr_{G}$ 可以等同于射影 $R\tl{t}$-子模 $\L\subset R\lr{t}^{n}$ 的模空间，其中对某个 $m\geq 0$ 有
\begin{equation}\label{lattice range}
(t^{m}R\tl{t})^{n}\subset \L\subset (t^{-m}R\tl{t})^{n}.
\end{equation}
这样的 $R\tl{t}$-模 $\L$ 称为 $R\lr{t}^{n}$ 中的{\em 格}。对固定的 $m$，令 $X_{m}$ 为 $\Gr_{G}$ 中分类满足 \eqref{lattice range} 的 $\L$ 的子函子，则 $X_{m}$ 可由 $k$ 上射影概形表示。自然的闭嵌入 $X_{m}\incl X_{m+1}$ 使 $\{X_{m}\}$ 成为射影概形的归纳系，而 $\Gr_{G}$ 可由归纳概形 $\varinjlim_{m}X_{m}$ 表示。

我们来详述 $\Gr_{G}(k)$ 与 $F^{n}=k\lr{t}^{n}$ 中的格之间的双射。设 $\calO_{F}^{n}\subset F^{n}$ 为标准格。设 $\frL_{n}$ 为 $F^{n}$ 中格的集合（当 $R=k$ 时，格就是 $F^{n}$ 的秩 $n$ 的 $\calO_{F}$-子模）。群 $LG(k)=G(F)$ 通过 $LG\ni g:\L\mapsto g\L$ 作用在 $\frL_{n}$ 上。此作用可迁，且标准格 $\calO^{n}_{F}$ 的稳定化子是 $L^{+}G(k)=G(\calO_{F})$。因此该作用诱导出 $G(F)$ 等变双射
\begin{equation}\label{Gr lattices}
\Gr_{G}(k)=G(F)/G(\calO_{F})\isom \frL_{n}.
\end{equation}

对一般的 $G$，$\Gr_{G}$ 总可由形如 $\varinjlim_{m}X_{m}$ 的归纳概形表示，其中 $X_{m}$ 是 $k$ 上射影概形，转移映射 $X_{m}\incl X_{m+1}$ 是闭嵌入。仿射 Schubert 分层给出 $\Gr_{G}$ 被射影概形的典范穷尽，现回忆如下。$L^{+}G$ 通过左平移作用在 $\Gr_{G}$ 上，其轨道以支配余权 $\l\in\xcoch(T)^{+}$ 为指标。用 $\Gr_{G,\l}$ 记过 $\l(t)$ 的 $L^{+}G$ 轨道，用 $\Gr_{G,\leq \l}$ 记 $\Gr_{G,\l}$ 的闭包。则 $\Gr_{G,\leq\l}$ 是射影概形，且 $\Gr_{G}$ 是这些 $\Gr_{G,\leq\l}$ 之并。关于仿射 Grassmann 簇的更多细节，参见 \cite[\S2]{BL}、\cite[\S2]{Faltings} 与 Zhu 的讲义 \cite{Zhu}。

\subsubsection{parahoric 子群} 子群 $L^{+}G\subset LG$ 是一类称为{\em parahoric 子群}的子群的例子。固定 Borel 子群 $B\subset G$，令 $\bI\subset L^{+}G$ 为 $B$ 在「模 $t$ 归约」映射 $L^{+}G\to G$ 下的原像。则 $\bI$ 是 $LG$ 的{\em Iwahori 子群}的一例。一般的 Iwahori 子群是 $\bI$ 在 $LG$ 中的共轭。与 $L^{+}G$ 一样，Iwahori 子群是 $LG$ 的无穷型群子概形。parahoric 子群是 $LG$ 中包含某个 Iwahori 子群、余维有限的连通群子概形。parahoric 子群的精确定义需要相当多的 Bruhat--Tits 理论，参见 Bruhat 与 Tits 的原始论文 \cite{BT1} 以及综述文章 \cite{Tits}。

正如 $G$ 的抛物子群的共轭类与 $G$ 的 Dynkin 图的子集一一对应，$LG$ 的 parahoric 子群在 $LG$ 共轭下的类，与 $G$ 的{\em 扩张 Dynkin 图} $\wt{\Dyn}(G)$（比 $G$ 的 Dynkin 图多一个顶点）的顶点的{\em 真子集}一一对应。扩张 Dynkin 图参见 Kac 的书 \cite[\S4.8]{Kac} 与 Bourbaki \cite[Ch VI]{Bourbaki}；与 parahoric 子群的联系参见 Gross 的综述文章 \cite{Gross}。

每个 $\bP$ 都有典范的群概形正合列
\begin{equation*}
1\to \bP^{+}\to \bP\to L_{\bP}\to1
\end{equation*}
其中 $\bP^{+}$ 是 $\bP$ 的{\em 顺幂幺根}，$L_{\bP}$ 是 $k$ 上的约化群，称为 $\bP$ 的{\em Levi 商}。若 $\bP$ 对应 $\wt{\Dyn}(G)$ 的顶点子集 $J$，则 Levi 商 $L_{\bP}$ 的 Dynkin 图是 $\wt{\Dyn}(G)$ 中由 $J$ 张成的子图。

\subsubsection{仿射旗簇} 对每个 parahoric 子群 $\bP\subset LG$，可以定义相应的{\em 仿射部分旗簇} $\Fl_{\bP}$，它是函子 $R\mapsto LG(R)/\bP(R)$ 在 $k$ 代数范畴中关于 fpqc 拓扑的层化。该函子同样可由归纳概形 $\varinjlim_{m}X_{m}$ 表示，其中每个 $X_{m}$ 是 $k$ 上射影概形，转移映射为闭嵌入。当 $\bP=L^{+}G$ 时，仿射 Grassmann 簇 $\Gr_{G}$ 是 $\Fl_{\bP}$ 的特例。

考虑特例：$\bP=\bI$ 是 $LG$ 的 Iwahori 子群。当 $G$ 单连通时，$\bI$ 是其自身的正规化子，于是可以把 $\Fl_{\bI}$ 等同于 $LG$ 的 Iwahori 子群的模空间，从而得到仿射旗簇的内蕴定义。通常把 $\Fl_{\bI}$ 记为 $\Fl$ 或 $\Fl_{G}$，称为{\em $G$ 的仿射旗簇}；但要提醒的是，只有当 $G$ 单连通时，$\Fl_{\bI}$ 才典范地不依赖于 $\bI$ 的选取。

设 $\bP\subset\bQ$ 是 $LG$ 的两个 parahoric 子群。则有自然投影 $\Fl_{\bP}\to \Fl_{\bQ}$。此投影的纤维同构于 $L_{\bQ}$ 的部分旗簇，对应的抛物子群由 $\bP\to L_{\bQ}$ 的像给出。特别地，有自然投影 $\Fl_{G}\to \Gr_{G}$，其纤维同构于旗簇 $\sB$。

\subsubsection{$\SL_{n}$ 的情形}\label{sss:GLn parahoric} 在 \S\ref{sss:Gr} 中已经看到，$\GL_{n}$ 的仿射 Grassmann 簇可以解释为格的模空间。事实上，$LG$ 的 parahoric 子群及相应的仿射部分旗簇也可以用格来描述。这里考虑 $G=\SL_{n}$ 的情形。

回忆 $F^{n}$ 中格的集合记为 $\frL_{n}$。对任意两个格 $\L_{1},\L_{2}\in\frL_{n}$，定义它们的相对长度为整数
\begin{equation*}
[\L_{1}:\L_{2}]:=\dim_{k}(\L_{1}/\L_{1}\cap\L_{2})-\dim_{k}(\L_{2}/\L_{1}\cap\L_{2}).
\end{equation*}
设 $J\subset \ZZ/n\ZZ$ 为非空子集。设 $\tilJ$ 是 $J$ 在投影 $\ZZ\to\ZZ/n\ZZ$ 下的原像。{\em 周期的 $J$-格链}指函数 $\L:\tilJ\to \frL_{n}$，把每个 $i\in \tilJ$ 送到格 $\L_{i}\in\frL_{n}$，并满足
\begin{itemize}
\item $[\L_{i}:\calO_{F}^{n}]=i$，对一切 $i\in \tilJ$；
\item $\L_{i}\subset \L_{j}$，对 $\tilJ$ 中的 $i<j$；
\item $\L_{i}=t\L_{i+n}$，对一切 $i\in \tilJ$。
\end{itemize}
设 $\frL_{J}$ 为周期的 $J$-格链的集合。对每个 $\{\L_{i}\}_{i\in \tilJ}\in\frL_{J}$，令 $\bP_{\{\L_{i}\}_{i\in J}}\subset LG$ 为所有 $\L_{i}$ 的同时稳定化子。则 $\bP_{\{\L_{i}\}_{i\in J}}$ 是 $LG$ 的 parahoric 子群，称之为 $J$ 型 parahoric 子群。事实上 $LG$ 的所有 parahoric 子群都来自唯一的一条周期 $J$-格链，其中 $J\subset \ZZ/n\ZZ$ 是非空子集且唯一。于是得到 $\sqcup_{J}\frL_{J}$ 与 $LG$ 的 parahoric 子群集合之间的双射。特别地，$L^{+}G$ 对应由 $\L_{i}=t^{i/n}\calO^{n}_{F}$（$i\in \tilJ=n\ZZ$）给出的周期 $\{0\}$-格链。

$G$ 的扩张 Dynkin 图是含 $n$ 个节点的圈，我们用 $\ZZ/n\ZZ$ 循环地为节点编号，使 $0$ 对应比通常 Dynkin 图多出的那个节点。$J\neq\varnothing$ 型的 parahoric 子群对应扩张 Dynkin 图节点的真子集 $\ZZ/n\ZZ-J$。

可以找到周期 $J$-格链的模空间 $\Fl_{J}$，使 $\Fl_{J}(k)=\frL_{J}$。固定任意 $J$ 型 parahoric 子群 $\bP$，$\Fl_{J}$ 可等同于仿射部分旗簇 $\Fl_{\bP}$。特别地，当 $G=\SL_{n}$ 时，仿射旗簇 $\Fl$ 可以等同于{\em 周期完备格链}的模空间，即 $F^{n}$ 中满足 $[\L_{i}:\calO^{n}_{F}]=i$ 与 $\L_{i}=t\L_{i+n}$（对一切 $i\in\ZZ$）的格序列 $\cdots \L_{-1}\subset \L_{0}\subset \L_{1}\cdots$。

\subsubsection{$\Sp_{2n}$ 的情形}\label{sss:Sp parahoric} 现考虑 $G=\Sp_{2n}=\Sp(V)$，其中 $V=k^{2n}$ 配有辛形式。把 $V$ 上的辛形式 $F$ 线性地扩张为 $V\otimes_{k}F$ 上的辛形式 $\jiao{-,-}$。对格 $\L\in \frL_{2n}$，定义其辛对偶为集合 $\L^{\vee}:=\{v\in V\otimes_{k}F|\jiao{v,\L}\subset \calO_{F}\}$。则 $\L^{\vee}$ 仍是 $V\otimes_{k}F$ 中的格。运算 $\L\mapsto \L^{\vee}$ 定义了 $\frL_{2n}$ 上的对合。

设 $J\subset\ZZ/2n\ZZ$ 是非空且在乘 $-1$ 下稳定的子集。设 $\tilJ\subset \ZZ$ 是 $J$ 在自然投影 $\ZZ\to\ZZ/2n\ZZ$ 下的原像。$V\otimes_{k}F$ 中的{\em 周期自对偶 $J$-格链}指满足下述附加条件的周期 $J$-格链（即按 \S\ref{sss:GLn parahoric} 的记号属于 $\frL_{J}$ 的元素）：
\begin{equation*}
\L^{\vee}_{i}=\L_{-i}, \text{ 对一切 } i\in \tilJ.
\end{equation*}
把 $V\otimes_{k}F$ 中周期自对偶 $J$-格链的集合记为 $\frL^{\Sp(V)}_{J}$。这是带 $G(F)=\Sp(V\otimes_{k}F)$ 作用的集合。对任意 $\{\L_{i}\}_{i\in\tilJ}\in \frL^{\Sp(V)}_{J}$，诸 $\L_{i}$ 的同时稳定化子是 $LG$ 的 parahoric 子群，而且 $LG$ 的每个 parahoric 子群都这样得到。对 $J$ 型 parahoric 子群 $\bP$，相应的仿射部分旗簇 $\Fl_{\bP}$ 可以等同于周期自对偶 $J$-格链的模空间，使 $\Fl_{\bP}(k)\cong\frL^{\Sp(V)}_{J}$ 作为 $G(F)$-集。请读者自行写出正交群的类似叙述，见习题 \ref{exII:SO parahoric}。
